\chapter*{TD4 : VARs, loi, moments.}

\setcounter{exercice}{0}

\begin{mdframed}
\begin{exercice} 
Posons, $n \in \mathbb{N}^*$ et $t > 0$ réel :

$$
I(n) = \int_t^\infty \frac{1}{(n - 1)!} x^{n-1}e^{-x} dx.
$$

\begin{enumerate}
    \item Montrer que l'on a :
    
    $$
    I(n) = \frac{1}{(n - 1)!} t^{n-1}e^{-t} + I(n - 1).
    $$
    \item En déduire :
    
    $$
    I(n) = \sum_{k=0}^{n-1} e^{-t} \frac{t^k}{k!}.
    $$
    \item Soit $X_t \sim \mathcal{P}(t)$ et $Y_n \sim \Gamma(n, 1)$. Montrer que :
    
    $$
    \mathbb{P}(X_t < n) = \mathbb{P}(Y_n > t).
    $$
\end{enumerate}
\end{exercice}
\end{mdframed}

\medskip

\begin{mdframed}
\begin{exercice}
\begin{enumerate}
    \item On cherche à caractériser l'ensemble des lois des variables aléatoires réelles $X$ qui satisfont à l'égalité en loi :
    
    $$
    X \overset{\text{loi}}{=} X^2.
    $$
    \begin{enumerate}
        \item Quel est le signe de $X$ ? En déduire la fonction de répartition $F$ de $X$ sur $]-\infty, 0[$.
        \item Montrer que $F$ satisfait $F(t) = F(t^2)$ pour tout $t \geq 0$.
        \item En déduire que $F$ est constante sur les deux intervalles $[0, 1[$ et sur $[1, \infty[$.
        \item Conclure : trouver l'ensemble des lois solutions.
    \end{enumerate}
    \item Adapter le raisonnement précédent pour les deux équations en loi suivantes :
    \begin{enumerate}
        \item $X \overset{\text{loi}}{=} X^3$
        \item \textit{(plus difficile, en bonus)} $X \overset{\text{loi}}{=} -X^3$
    \end{enumerate}
\end{enumerate}
\end{exercice}
\end{mdframed}

\medskip

\begin{mdframed}
\begin{exercice}
On se donne $X_1, \dots, X_n$, $n$ variables aléatoires indépendantes de loi $\mathcal{U}(0, 1)$ et on note $X_{(1)} \leq \dots \leq X_{(n)}$ leur réordonnancement croissant (encore appelé statistique d'ordre) ; ainsi, $X_{(1)}$ est le plus petit des $X_i$, $i = 1 \dots n$, $X_{(2)}$ le deuxième plus petit...

\begin{enumerate}
    \item Montrer que $\mathbb{P}(X_{(n)} \leq t) = t^n$ pour tout $t \in [0, 1]$ et en déduire que $X_{(n)}$ suit une loi Bêta dont on précisera les paramètres.
    \item Montrer que $\mathbb{P}(X_{(n-1)} \leq t) = t^n + n(1 - t)t^{n-1}$ pour tout $t \in [0, 1]$ et en déduire que $X_{(n-1)}$ suit une loi Bêta dont on précisera les paramètres.
    \item De même, donner les lois de $X_{(1)}$ et $X_{(2)}$.
    \item Quel résultat peut-on conjecturer de façon plus générale au sujet de la loi de $X_{(i)}$ si $i \in \{1, \dots, n\}$ ?
\end{enumerate}
\end{exercice}
\end{mdframed}

\medskip

\begin{mdframed}
\begin{exercice}
Soit $\alpha, \beta > 0$ et $n \in \mathbb{N}^*$. On rappelle la définition de la fonction Gamma d'Euler :

$$
\Gamma(t) = \int_0^\infty x^{t-1}e^{-x} dx, \quad t > 0.
$$

\begin{enumerate}
    \item Montrer que $\Gamma(n + 1) = n!$.
    \item Calculer le $n$-ième moment de la loi exponentielle $\mathcal{E}(\beta)$, c'est-à-dire $\mathbb{E}[X^n]$ si $X \sim \mathcal{E}(\beta)$.
    \item Calculer le $n$-ième moment de la loi Gamma $\Gamma(\alpha, \beta)$.
    \item Comparez les résultats obtenus aux deux dernières questions.
\end{enumerate}
\end{exercice}
\end{mdframed}

\medskip

\begin{mdframed}
\begin{exercice}
\begin{enumerate}
    \item Calculer les moments impairs de la loi normale centrée réduite $\mathcal{N}(0, 1)$.
    \item On cherche maintenant à calculer les moments pairs de $\mathcal{N}(0, 1)$.
    \begin{enumerate}
        \item Montrer à l'aide d'une intégration par parties que :
        
        $$
        \int_{\mathbb{R}} x^{2n}e^{-\frac{x^2}{2}} dx = (2n - 1) \int_{\mathbb{R}} x^{2n-2}e^{-\frac{x^2}{2}} dx
        $$
        \item En déduire que le $2n$-ième moment de $\mathcal{N}(0, 1)$ vaut :
        
        $$
        \prod_{k=1}^n (2k - 1) = \frac{(2n)!}{2^n \cdot n!}
        $$
    \end{enumerate}
\end{enumerate}
\end{exercice}
\end{mdframed}